% Generated draft: requires mathematical and source review.
\documentclass[11pt]{article}
\usepackage[margin=25mm]{geometry}
\usepackage{fontspec}
\setmainfont{Latin Modern Roman}
\usepackage{amsmath,amssymb,mathtools,unicode-math}
\setmathfont{Latin Modern Math}
\usepackage{enumitem,needspace,xurl,hyperref}
\setlength{\parindent}{0pt}
\setlength{\parskip}{5pt}
\setlength{\emergencystretch}{3em}
\newcommand{\N}{\mathbb{N}}
\newcommand{\Z}{\mathbb{Z}}
\newcommand{\Q}{\mathbb{Q}}
\newcommand{\R}{\mathbb{R}}
\newcommand{\C}{\mathbb{C}}
\newcommand{\F}{\mathbb{F}}
\newcommand{\A}{\mathbb{A}}
\newcommand{\PP}{\mathbb{P}}
\newcommand{\E}{\mathbb{E}}
\newcommand{\eps}{\varepsilon}
\newcommand{\dd}{\,\mathrm{d}}
\begin{document}
\section*{UnsolvedMath ID: 30001577}
\textbf{Source problem code:} catalogue entry
\subsection{Definitions and mathematical statement}
Let \(k\) be a field of characteristic \(p>0\), where \(p\) is prime; the record does not assume that \(k\) is finite. Let \(\Gamma\) be a residually finite group, meaning that every nonidentity element survives in some finite quotient of \(\Gamma\). Consider a descending sequence
\[
\Gamma=\Gamma_0\geq\Gamma_1\geq\Gamma_2\geq\cdots,
\qquad [\Gamma:\Gamma_i]<\infty,
\qquad \bigcap_{i\geq0}\Gamma_i=\{1\}.
\]
The supplied statement calls this a residual chain but does not explicitly specify whether every \(\Gamma_i\) must be normal in \(\Gamma\). Thus the admissible class of chains, and consequently the scope of the independence question below, requires source confirmation.

Let \(X\) be a finite free \(\Gamma\)-CW complex. Under the usual equivariant interpretation of this terminology, the action is free and there are finitely many \(\Gamma\)-orbits of cells, so that \(\Gamma\backslash X\) is a finite CW complex. This interpretation of ``finite'' requires confirmation from the source; finiteness of the underlying CW complex is a different condition. No contractibility or connectedness hypothesis on \(X\) is supplied.

For each integer \(n\geq0\), put
\[
b_n(\Gamma_i\backslash X;k)
:=\dim_k H_n(\Gamma_i\backslash X;k),
\]
where \(H_n(-;k)\) denotes ordinary singular homology with coefficients in \(k\), equivalently cellular homology for these CW complexes. The index \([\Gamma:\Gamma_i]\) is the number of cosets of \(\Gamma_i\) in \(\Gamma\).

The supplied question asks, for every such \(k\), \(\Gamma\), and \(X\), and every integer \(n\geq0\):
\begin{enumerate}
\item Does the finite real limit
\[
L_n(\Gamma,X;k;(\Gamma_i))
:=\lim_{i\to\infty}
\frac{b_n(\Gamma_i\backslash X;k)}{[\Gamma:\Gamma_i]}
\]
exist for every admissible residual chain?
\item Whenever these limits exist, are they independent of the residual chain? More explicitly, for two admissible residual chains \((\Gamma_i)\) and \((\Lambda_j)\) in the same group, must
\[
\lim_{i\to\infty}\frac{b_n(\Gamma_i\backslash X;k)}{[\Gamma:\Gamma_i]}
=
\lim_{j\to\infty}\frac{b_n(\Lambda_j\backslash X;k)}{[\Gamma:\Lambda_j]}?
\]
\item If a chain-independent limit exists, what invariant of the data \((\Gamma,X,k,n)\) equals it?
\end{enumerate}
The record supplies no proposed invariant or precise candidate equality. Its original wording asks about ``every group,'' whereas its normalized statement restricts to residually finite groups equipped with a residual chain. No formulation for groups without such a chain is supplied, so an extension to those groups remains unspecified.
\subsection{Short English statement}
For a field of positive characteristic and a free group action with finite CW quotient, ask whether the Betti numbers of the finite covers, divided by their covering degrees, converge along every residual chain. Ask also whether the limit is independent of the chain and which invariant describes it. The intended normality condition on the chain and the finiteness convention require confirmation. The supplied open label is dated curation, not an independently verified status.
\subsection{Sources}
\begin{itemize}
\item Supplied local UnsolvedMath registry record: permanent ID \(30001577\), title ``Characteristic-p Approximation of Normalized Betti Numbers,'' exact source problem code \texttt{catalogue entry}, in the Oberwolfach Reports Open Problems collection. No individual proposer is supplied; the recorded proposed year is \(2010\).
\item Supplied report citation and literature source labeled ``original OWR report'': \emph{Actions and Invariants of Residually Finite Groups: Asymptotic Methods} (2011), Oberwolfach Reports, pp.~2335--2391, DOI \texttt{10.4171/owr/2010/40}. \url{https://doi.org/10.4171/owr/2010/40}. The record attributes the problem to workshop open-problem discussions but gives no precise problem page or speaker attribution. The supplied citation year \(2011\), proposed year \(2010\), and DOI containing \(2010\) are retained without reconciliation.
\item Supplied literature source labeled ``Lueck, approximation of L2-invariants'': \url{https://arxiv.org/abs/math/0111039}. No exact theorem or passage supporting a positive-characteristic formulation is identified in the supplied record.
\item Literature assessment checked on \texttt{2026-08-21}: the supplied curation labels the problem open and describes the general arbitrary-chain limit problem as unsettled, while referring to Lueck--Farber approximation in important Farber or free-action settings. This is supplied dated curation, not independently verified; the applicability and coefficient hypotheses of those approximation results require independent checking.
\item The supplied statement status is \texttt{corrected\_verified}; its verification note reports that the mathematical content was checked against the official Oberwolfach report and its notation, context, and question form normalized. This is a supplied verification claim, not verification performed for this draft. The background explicitly requests independent verification of time-sensitive status and cited claims.
\end{itemize}
\end{document}
